\chapter{Conclusion and Future Work}

\section{Conclusion}
In this work, we present a physics model capturing three key aspects of a 
two-phase partially molten mantle: (1) a degenerate Darcy–Stokes system 
describing the static mechanical behavior, (2) a nonlinear  
system governing the transport of chemical species and thermal energy, 
and (3) a eutectic phase relation characterizing melting processes.
We develop a local mass conservative numerical framework to address this 
complicated multiphysics system. 

For the static mechanics system, we introduce an $H(\text{div})$-and an $H^1$- 
conforming mixed finite element spaces for Darcy and Stokes subproblems 
respectively. The mixed finite element spaces are constructed using a noval 
direct serendipity space tailored for quadrilateral mesh.  
A modified Bernardi–Raugel element is developed and analyzed. 
We provide detailed proofs of its stability and approximation properties.
A new modified uzawa iterative method is implemented to solve the 
singular saddle 
point problem induced by this degenerate coupled system. 

To address the nonlinear transport of chemical species and energy in the 
presence of porous media, we study and modify the weighted essentially 
non-oscillatory (WENO) method.
We propose two simple alternatives to accelerate the evaluation of the 
smoothness indicator. Additionally, we develop a new multilevel WENO weighting 
scheme that is more versatile and flexible, allowing the combination of 
arbitrary stencils of different orders.

We advance the multiphysics system by incorporating a eutectic phase module that
 couples compositional and thermal information with porosity, thereby capturing 
the evolution of porosity and phase transitions.

The code is implemented in C/C++ along with PETSc, 
and is designed to be easily adapted 
for large-scale parallel execution on supercomputers,
 with the potential to support general unstructured meshes.

\section{Future Work}

We propose to use the new framework to sovle some important problem, 
targeting applicatoins in simulation of partial melting materials, not 
partial molten mantle but glacier as well.

We also propose to extend our framework to three-dimensional and 
general unstructured mesh.
